\Article{25}
Mit \(\lambda,\lambda',\lambda''\) etc. bezeichnen wir ferner die Werte der Funktionen \(v,v',v''\) etc., die sich ergeben, wenn für \(x,y,z\) etc. ihre plausibelsten Werte \(A,B,C\) etc. eingesetzt werden, nämlich
\SourceLeaf{037}
\[
\begin{aligned}
aA+bB+cC+\text{etc.}+l&=\lambda\\
a'A+b'B+c'C+\text{etc.}+l'&=\lambda'\\
a''A+b''B+c''C+\text{etc.}+l''&=\lambda''
\end{aligned}
\]
etc.; setzen wir außerdem
\[
\lambda\lambda+\lambda'\lambda'+\lambda''\lambda''+\text{etc.}=M
\]
sodass \(M\) der Wert der Funktion \(\Omega\) ist, der den plausibelsten Werten der Unbestimmten entspricht, und somit nach dem, was wir in Art. 20 bewiesen haben, der kleinste Wert dieser Funktion. Dann wird \(a\lambda+a'\lambda'+a''\lambda''+\text{etc.}\) der Wert von \(\xi\) sein, der den Werten \(x=A,\allowbreak y=B,\allowbreak z=C\) etc. entspricht, und folglich \(=0\); wir haben also
\[
\Sigma a\lambda=0
\]
und ebenso gilt
\[
\begin{aligned}
&\Sigma b\lambda=0,\quad\Sigma c\lambda=0\quad\text{etc.;}\\
&\text{sowie}\quad\Sigma\alpha\lambda=0,\quad\Sigma\beta\lambda=0,\quad\Sigma\gamma\lambda=0\quad\text{etc.}
\end{aligned}
\]
Multiplizieren wir schließlich die Ausdrücke für \(\lambda,\lambda',\lambda''\) etc. jeweils mit \(\lambda,\lambda',\lambda''\) etc. und addieren, so erhalten wir \(l\lambda+l'\lambda'+l''\lambda''+\text{etc.}=\lambda\lambda+\lambda'\lambda'+\lambda''\lambda''+\text{etc.}\) oder
\[
\Sigma l\lambda=M
\]
\Article{26}
Setzen wir in der Gleichung \(v=ax+by+cz+\text{etc.}+l\) für \(x,y,z\) etc. die Ausdrücke VII aus Art. 21 ein und wenden die aus dem Vorangehenden offensichtlichen Reduktionen an, so ergibt sich
\[
v=\alpha\xi+\beta\eta+\gamma\zeta+\text{etc.}+\lambda
\]
und ebenso gilt identisch
\[
\begin{aligned}
v'&=\alpha'\xi+\beta'\eta+\gamma'\zeta+\text{etc.}+\lambda'\\
v''&=\alpha''\xi+\beta''\eta+\gamma''\zeta+\text{etc.}+\lambda''
\end{aligned}
\]
etc. Multiplizieren wir entweder diese Gleichungen oder die Gleichungen I aus Art. 20 jeweils mit \(\lambda,\lambda',\lambda''\) etc. und addieren, so finden wir die Identität
\[
\lambda v+\lambda'v'+\lambda''v''+\text{etc.}=M
\]
\clearpage
